% Conversão fiel da participação como cliente Hoortech; origem imutável RELATORIO_ATUAL.odt.
\hypertarget{__RefHeading___Toc4470_311817663321}{}

% SOURCE text-h-0029
\refstepcounter{subsection}
\subsection*{6.2.3 Sprint 1}
\label{cliente-hoortech:text-h-0029}
\addcontentsline{toc}{subsection}{6.2.3 Sprint 1}

% SOURCE text-p-0718
\originaltab{}Na Sprint 1, os alunos avançaram significativamente, com foco no desenvolvimento e aprimoramento do front-end. Praticamente todos os elementos dessa interface foram finalizados, embora alguns ajustes técnicos, como o funcionamento da câmera para captura dos sinais ainda estivessem sendo refinados. Importante ressaltar que todas as sugestões de melhoria apresentadas por mim e outros colegas durante a entrega anterior foram atendidas prontamente. Esse fato demonstrou que os alunos ouviram nosso feedback e implementaram as mudanças necessárias, mostrando proatividade e compromisso com o projeto.\originalspace{}\par

% SOURCE text-p-0719
\originaltab{}Em relação ao back-end, a equipe deu início à construção dessa camada da aplicação, que será responsável pela integração e suporte ao modelo de reconhecimento de linguagem de sinais brasileiras (Libras). Embora o desenvolvimento do back-end ainda estivesse em uma fase inicial e, portanto, sem elementos visuais ou integrações para demonstrar, o progresso notável no front-end trouxe uma grande tranquilidade e confiança no desenvolvimento do projeto como um todo.\par

